\section{The Dirichlet problem on a periodic fiber}
\label{sec:dirichlet}

The balance law in Section~\ref{sec:energy} uses a Dirichlet potential
smooth up to a boundary that moves with the time variables.
Section~\ref{sec:barrier} bounds this potential by comparing
determinants on the whole fiber and on small balls inside it. We prove
the required existence, comparison, and smooth dependence here;
the quantitative bounds in the wavelength $h$ come from the
comparison functions in Section~\ref{sec:barrier}. Alesker proved smooth
Dirichlet solvability on quaternionic balls \cite[Theorem~1.4]{Alesker2003MA};
Zhu proved it on bounded domains in quaternionic space in the presence of
a smooth subsolution \cite[Theorem~1.1]{Zhu2016}. Our fibers have periodic
real directions and only three real coordinates per variable. Adjoining a
fourth, periodic coordinate gives the local rotations needed to estimate
mixed derivatives at the boundary.

Let $\Lambda\subset\RR^m$ be a full lattice, $V\subset\RR^{2m}$ a
bounded domain with smooth boundary, and
$\Omega=(\RR^m/\Lambda)\times V$. The periodic variables occupy the
real slots and the coordinates of $V$ occupy the $\mathbf i,\mathbf j$
slots of the Hessian $H$ from Lemma~\ref{lem:hessian-calculus}.

\begin{theorem}[Periodic Dirichlet problem]\label{thm:dirichlet}
Let $m\ge1$. Suppose $r\in C^\infty(\overline V)$ is a defining
function, negative in $V$, with $dr\ne0$ on $\partial V$ and
$H(r)>0$ on $\overline\Omega$. Let $B_0,g\in C^\infty(\overline\Omega)$
be real and suppose
\[
 H(B_0)\ge0,\qquad g>\det H(B_0)\quad\text{on }\overline\Omega.
\]
There is a unique smooth solution, up to the boundary, of
\begin{equation}\label{eq:dirichlet-problem}
 \det H(B_0+u)=g,\qquad H(B_0+u)>0,\qquad
 u|_{\partial\Omega}=0.
\end{equation}
There are constants $b\ge b'>0$ such that
\begin{equation}\label{eq:dirichlet-barriers}
 br\le u\le b'r.
\end{equation}
For smooth finite-dimensional families of such data and domains, these
constants can be chosen locally uniformly. The solutions depend smoothly
on the parameters up to the boundary, after any smooth local
identification of the domains with a fixed domain.
\end{theorem}

\begin{proof}
Work on $\widetilde\Omega=\Omega\times(\RR/\mathbb Z)^m$, with the new
coordinates in the $\mathbf k$ slots. Extend the data constantly in these
coordinates. Until the last paragraph, $G,H$ are the full quaternionic
gradient and Hessian, which agree with the physical operators on the
extended data. Write $\mathcal H(D^2v)=H(v)$ and encode a real covector
$\xi\in\RR^{4m}$ as the quaternionic column $\widehat\xi$ in the slot
basis. Lemma~\ref{lem:hessian-calculus} gives
$\mathcal H(\xi\xi^{\mathsf t})=\widehat\xi\,\widehat\xi^*$ and
$\tr H(v)=\Delta v$.

\emph{Comparison and continuity.}
Join the known solution $u=r$ to the given density by
$g_t=(1-t)\det H(B_0+r)+tg$, $0\le t\le1$.
At a solution $\psi=B_0+u$ put $Q=H(\psi)$ and $S=\tr Q^{-1}$.
The derivative of the log determinant is
\[
 Lw=Q^{-1}:H(w),\qquad
 \sigma_L(\xi)=\widehat\xi^*Q^{-1}\widehat\xi .
\]
Thus $L$ is a uniformly elliptic real operator at each fixed smooth
solution with $Q>0$ on the closure. For two potentials $\psi_1,\psi_2$
with $Q_i=H(\psi_i)>0$, the corresponding difference formula is
\[
 \log\det Q_1-\log\det Q_2
   =\left(\int_0^1\big((1-\theta)Q_2+\theta Q_1\big)^{-1}\dd\theta\right)
          :H(\psi_1-\psi_2).
\]
When the Hessians are positive on the closure, the averaged operator
is uniformly elliptic. If $\det Q_1\ge\det Q_2$ on the domain and
$\psi_1\le\psi_2$ on its boundary, the maximum principle therefore
gives $\psi_1\le\psi_2$ throughout. This also gives uniqueness.
The linear Dirichlet
operator $L:C^{k+2,\alpha}_0\to C^{k,\alpha}$ is an isomorphism
($k\ge0$, $0<\alpha<1$; the subscript denotes zero boundary trace).
Indeed the local Schauder estimates and Dirichlet theory
\cite[Theorems~6.14 and~6.19]{GilbargTrudinger} apply in a finite
cover of the flat quotient. In the continuity method from $\Delta$ to
$L$, $r$ is a common strict barrier since $H(r)\ge cI$; the maximum
principle therefore supplies the $C^0$ estimate. The starting Dirichlet
Laplacian is invertible on the compact manifold with boundary.
The implicit function theorem gives openness; the estimates below give
closedness. Their constants may depend on the fiber and its data, but
are uniform in $t$.

Both endpoints of the path have determinant strictly above $\det H(B_0)$
at every point. The gap is uniform by compactness. Hence we may choose
$0<b'<1<b$ so that
\[
 \det H(B_0+b'r)<g_t<\det H(B_0+br)
\]
everywhere, using $H(r)\ge cI$ for the second inequality. Comparison
gives \eqref{eq:dirichlet-barriers}. Since $u$ vanishes on the boundary,
$u=\alpha r$ in a collar with smooth quotient $\alpha$ at each smooth
solution; division in the barriers gives $b'\le\alpha\le b$. In
particular $\nabla u=\alpha\nabla r$ on the boundary.
We also have $Lr\ge cS$ and $S\ge m g_t^{-1/m}\ge c_0>0$.
For any constant real unit vector $e$, differentiation and concavity give
\begin{equation}\label{eq:dirichlet-directional}
 L\psi_e=(\log g_t)_e,\qquad L\psi_{ee}\ge(\log g_t)_{ee}.
\end{equation}
The maximum principle applied to $\pm\psi_e+Cr$ bounds the full
gradient. Applied to $\psi_{ee}+Cr$, it reduces an upper bound for the
real Hessian to the boundary.

\emph{Second derivatives at the boundary.}
Center a flat chart at a boundary point, let $\nu$ be the outward unit
normal, and identify vectors with their quaternionic columns in this
chart. For real tangent vectors $v,w$,
$D^2u(v,w)=\alpha D^2r(v,w)$ at the center. For the mixed derivatives
we adapt the tangential differentiation argument of Caffarelli, Nirenberg
and Spruck for the real Monge--Ampère equation
\cite[Section~2, (2.11)--(2.14)]{CNS1984}. Fix a real tangent unit
vector $e$. We construct a local quaternionic rigid motion
$X=(e+Ax)\cdot\nabla$ for which $Xr$ vanishes to second order along
the boundary at the center. Let
\[
 v=\frac{P_{\rm tan}D^2r\,e}{|\nabla r|},\qquad
 A=v\nu^*-\nu v^*-\nu(\nu^*v)\nu^* ,
\]
where $P_{\rm tan}$ is the real orthogonal tangent projection. Since
$\Rea(\nu^*v)=0$, we have $A^*=-A$ and $A\nu=v$. Now $Xr(0)=0$ and
for every tangent $w$,
\[
 \partial_w(Xr)(0)=D^2r(w,e)-|\nabla r|\,w\cdot A\nu=0.
\]
Thus $|Xu|=|\alpha Xr|\le C|x|^2$ on the boundary near the center.
The flow of $X$ acts by quaternionic unitary rotations and translations.
Differentiating the covariance of $H$ under this flow gives
$L(X\psi)=X\log g_t$: the commutator from conjugating $Q$ has zero
pairing with $Q^{-1}$. This identity uses only the flat chart.
Subtracting the known $XB_0$ gives $|L(Xu)|\le CS$.

In a chart ball of fixed small radius, set $w=C_1|x|^2-C_2r$.
First choose $C_1$ so that $w\ge|Xu|$ on the true boundary by the
quadratic bound, and on the spherical boundary by the gradient bound.
Then choose $C_2$ so that $L(w\pm Xu)\le0$, using $H(|x|^2)=8I$
and $Lr\ge cS$. The maximum principle gives $|Xu|\le w$ in the
intersection. Both functions vanish at the center, so the inward
derivative bounds $|\partial_\nu Xu|$. This is $|D^2u(\nu,e)|$:
the additional term $\nabla u\cdot A\nu$ vanishes since $\nabla u$
is normal. The constants are uniform over the compact boundary.

For the remaining derivative, the product rule at the boundary reads
\[
 H(u)=\alpha H(r)+G(\alpha)G(r)^*+G(r)G(\alpha)^*,\qquad
 G(r)=|\nabla r|\nu .
\]
Compression to the \emph{quaternionic} orthogonal complement of $\nu$
annihilates both cross terms. On that complement $Q$ is therefore bounded
below by $b'H(r)|_{\nu^\perp}\ge cI$. The rank-one rule shows that in a
quaternionic orthonormal basis beginning with $\nu$,
\[
 Q=\begin{pmatrix}\psi_{\nu\nu}+c_1&z^*\\z&D\end{pmatrix}.
\]
Here $c_1,z,D$ depend only on the real second derivatives already
bounded. For $m>1$ the Schur formula in
Lemma~\ref{lem:quaternionic-matrices} gives
$\psi_{\nu\nu}=g_t/\det D+z^*D^{-1}z-c_1$, which is bounded since
$D\ge cI$. For $m=1$ the scalar identity $Q=g_t$ suffices.
Equation~\eqref{eq:dirichlet-directional} now gives $D^2\psi\le CI$
throughout the domain. Since $\Delta\psi=\tr Q>0$, every real Hessian
eigenvalue is at least $-(4m-1)C$. This is the full $C^2$ bound.

\emph{Real elliptic regularity.}
The $C^2$ bound and $\det Q=g_t\ge c_0>0$ put all eigenvalues of $Q$
in a fixed interval $[c,C]\subset(0,\infty)$. Thus all real Hessians
$D^2\psi$ lie in the compact convex set
$\mathcal K=\{M:|M|\le C_1,\ cI\le\mathcal H(M)\le CI\}$.
On $\mathcal H(M)>0$ the smooth function $F(M)=\log\det\mathcal H(M)$
is concave. For $M\in\mathcal K$ and every real covector $\xi$,
\[
 DF(M)[\xi\xi^{\mathsf t}]
   =\widehat\xi^*\mathcal H(M)^{-1}\widehat\xi
   \in[C^{-1}|\xi|^2,c^{-1}|\xi|^2].
\]
These are real uniform ellipticity and concavity on the range of the
solution Hessians, the hypotheses (i)$'$, (ii)$'$ of the interior and
boundary estimates \cite[Theorems~17.14 and~$17.26'$]{GilbargTrudinger}. To put
$F$ on the full real matrix space as in that theorem, multiply it by a
smooth cutoff supported in $\{\mathcal H(M)>0\}$ and equal to one on
a neighborhood of $\mathcal K$. The equation
$F(D^2\psi)-\log g_t=0$ then has uniformly smooth coefficients and
boundary data $B_0$. The cited local boundary estimates (see also
\citet{Krylov1983}) and the interior Evans--Krylov estimate
\cite{Evans1982,Krylov1982} apply in a finite
cover of the flat quotient, giving a uniform $C^{2,\alpha}$ bound for
some $\alpha>0$.

For a limit of solvable parameters, a subsequence converges in
$C^{2,\beta}$, $0<\beta<\alpha$. Positivity and the equation pass to
the limit. Schauder bootstrapping gives smoothness up to the boundary:
in flattened coordinates differentiate tangentially with the given
Dirichlet data, and recover the pure normal derivative from the equation
by uniform ellipticity. Thus the path is closed and reaches $g$.

Uniqueness on $\widetilde\Omega$ makes the solution invariant under
translations in every added periodic coordinate. It descends to
$\Omega$; any physical solution lifts, so uniqueness descends as well.
The same lift gives comparison on subdomains, in particular on the
balls in the three-slot fibers used in Section~\ref{sec:barrier}.
For moving domains pull the lifted equation back to a fixed domain.
Its coefficients and boundary data depend smoothly on the parameters,
and its linearization is conjugate to the Dirichlet isomorphism above.
The implicit function theorem in each $C^{k+2,\alpha}_0$ gives smooth
dependence up to the boundary. The strict inequalities used to choose
$b,b'$ persist locally in the parameters.
\end{proof}
